\documentclass[11pt]{article}
\usepackage{amsmath,amssymb,amsthm}
\newtheorem{theorem}{Theorem}
\begin{document}
\title{Generic stability of the zero-residue reciprocal branch}
\author{}\date{}\maketitle
\paragraph{Status and scope.}
This is a complete written, elementary argument, not a Lean theorem.
It concerns fixed query offsets and a sufficiently small neighborhood
of a zero-residue branch. It does not bound the original model, the
full affine-numerator family, or deliberately singular or scale-dependent
query configurations.

The common-zero identity
\[
 H+\sum_{h=1}^H\frac{a_h-D_{hj}}{D_{hj}-p_h}
 =\sum_{h=1}^H\frac{a_h-p_h}{D_{hj}-p_h}
\]
shows that every pole vector $p$ is a zero branch when $a=p$.
Writing $a=p+u$, its exact score space is the column space of
$V(p)_{jh}=1/(D_{hj}-p_h)$.

\begin{theorem}
Let $q\geq H\geq1$, let every $D_{hj}>0$, and suppose every
$H$-row minor of $V(0)$ is nonzero. There is $\delta>0$ such that
for every $p\in(-\delta,\delta)^H$, the set
\[
 \{\operatorname{sign}(V(p)u):u\in\mathbb R^H,
                  \ (V(p)u)_j\ne0\text{ for every }j\}
\]
is independent of $p$. In particular, allowing all such zero branches
adds no strict sign words to a single $H$-dimensional linear score space.
The number of words is
\[
 2\sum_{i=0}^{H-1}\binom{q-1}{i}.
\]
\end{theorem}
\begin{proof}
Choose $\delta$ so all denominators remain positive and every
$H$-row determinant retains its nonzero sign. This is possible by
continuity and finiteness. All subsets of at most $H$ rows remain
linearly independent: otherwise extending such a subset to $H$ rows
would contradict a nonzero $H$-row determinant.

For a proposed word $\sigma\in\{-1,1\}^q$, the inequalities
$\sigma_jV(p)_{j,\cdot}u>0$ are feasible exactly when the origin is
not in the convex hull of the signed rows. One direction follows by
taking the inner product with $u$; for the other, the closest point
of that compact convex hull to the origin gives a strictly separating
linear functional. If the origin belongs to the convex hull, choose
a convex representation with minimal support. Its rows are minimally
linearly dependent, since a second independent dependence could be
added until a positive coefficient becomes zero. Its support has
exactly $H+1$ elements: smaller supports are independent, and a
minimal dependence in $\mathbb R^H$ has at most $H+1$ elements.
For an ordered support $j_0,\ldots,j_H$, its unique dependence has
coefficients
\[
 (-1)^k\det V(p)_{\{j_0,\ldots,j_H\}\setminus\{j_k\},\cdot},
 \qquad 0\leq k\leq H,
\]
up to one common nonzero factor. Signing the rows multiplies these
coefficients by the prescribed signs, up to another common factor.
All determinant signs are unchanged, so the existence of a positive
dependence is unchanged. This proves equality of the strict word sets.
For $q=H$ all words occur directly by invertibility.

Finally, $q$ central hyperplanes in general position in $\mathbb R^H$
have $2\sum_{i=0}^{H-1}\binom{q-1}{i}$ regions. This follows by
adding the last hyperplane and applying the same recurrence to its
induced arrangement, with the one-dimensional and single-hyperplane
base cases. Each strict word is one convex open region.
\end{proof}

\paragraph{Genericity and the two-head collapse.}
The hypothesis holds on an open dense subset of positive query-offset
matrices. After clearing positive denominators, each determinant is a
polynomial which is not identically zero: choosing
$D_{hj}=x_j+y_h$ with distinct positive $x_j$ and distinct positive
$y_h$ gives a nonsingular Cauchy minor for every row subset. A finite
union of proper polynomial zero sets has empty interior.
For $H=2$ there is a stronger global feature identity, with no
small-neighborhood assumption:
\[
 \operatorname{sign}\left(\frac{u_1}{D_{1j}-p_1}
              +\frac{u_2}{D_{2j}-p_2}\right)
 =\operatorname{sign}\bigl(u_1D_{2j}+u_2D_{1j}
                      -u_1p_2-u_2p_1\bigr)
\]
whenever the two denominators are positive. Thus that subfamily has
VC dimension at most three. Neither statement covers nonzero fixed
column numerator offsets, whose full two-head family can shatter five
queries as proved elsewhere in this repository.

\paragraph{Exact check.}
The companion checker uses rational interval arithmetic to certify
that all $35$ three-row determinants of one $7\times3$ example retain
their signs throughout an explicit pole box. It also enumerates the
circuit criterion and obtains the predicted $44$ strict words.
The unresolved escape is a deliberately singular or scale-dependent
query arrangement: its determinant signs can change at arbitrarily
small pole scales. No claim of universality is made for that escape.
\end{document}
